\newcommand{\DoNotLoadEpstopdf}{}
\pdfinfoomitdate=1
\pdftrailerid{}
\pdfsuppressptexinfo=-1
\documentclass[11pt,letterpaper]{article}
\usepackage[T1]{fontenc}
\usepackage{lmodern}
\usepackage{amsmath,amssymb,amsthm,mathtools}
\usepackage[margin=1in]{geometry}
\usepackage{booktabs,longtable,array}
\usepackage{microtype}
\usepackage[hidelinks]{hyperref}
\hypersetup{pdftitle={Report 153 All orders refinement for unlabeled permutation graphs},pdfauthor={},pdfsubject={OEIS A123448 factorial expansion, inverse thresholds, and sharp total variation rate}}
\setlength{\parskip}{4pt}
\setlength{\parindent}{1.2em}
\setlength{\emergencystretch}{2em}
\setcounter{tocdepth}{1}
\numberwithin{equation}{section}
\newtheorem{theorem}{Theorem}[section]
\newtheorem{lemma}[theorem]{Lemma}
\newtheorem{proposition}[theorem]{Proposition}
\newtheorem{corollary}[theorem]{Corollary}
\theoremstyle{definition}\newtheorem{definition}[theorem]{Definition}
\theoremstyle{remark}\newtheorem{remark}[theorem]{Remark}
\newcommand{\PP}{\mathbb P}
\newcommand{\EE}{\mathbb E}
\newcommand{\ZZ}{\mathbb Z}
\newcommand{\QQ}{\mathbb Q}
\newcommand{\eps}{\varepsilon}
\newcommand{\ceil}[1]{\left\lceil #1\right\rceil}
\newcommand{\floor}[1]{\left\lfloor #1\right\rfloor}
\DeclareMathOperator{\Aut}{Aut}
\DeclareMathOperator{\TV}{TV}
\DeclareMathOperator{\Inv}{Inv}
\title{All orders refinement for unlabeled permutation graphs\\\large Factorial expansion, inverse thresholds, and sampling in OEIS A123448}
\author{Report 153}
\date{3 October 2026}
\begin{document}
\maketitle
\begin{abstract}
Let $a_n$ count unlabeled permutation graphs on $n$ vertices. We establish,
for every fixed $R\geq1$,
\[
 a_n=\frac14\sum_{k=0}^{R-1}h_k(n-k)!+O_R((n-R)!),
\]
with a formal coefficient generator using only finite unrooted and
vertex-rooted graph counts at each order. The first coefficients are
$1$, $-4$, $-1$, $-94/3$, $-769/6$, $-19969/15$,
$-531812/45$, and $-40114096/315$.
The main step is a finite-polynomial approximation: all but
$O_R((n-R)!)$ permutation realizers have a recoverable large prime
quotient, bounded decorations, and a bounded vertex-rooted context.
Borinsky's established factorial composition theorem then transfers
all fixed orders without assuming an asymptotic expansion of the
unknown graph series. We derive denominator divisibility,
arbitrary fixed-order threshold brackets, an explicit second
Lambert-$W$ inverse shift, and the sharp all-orders sampling comparison beginning
$d_{\TV}=4/n-15/n^2+O(n^{-3})$ between uniform graph classes and graphs of
uniform permutations. Exact companion programs provide independent
coefficient routes, formal inverse checks, and two exhaustive finite
enumeration options.
\end{abstract}
\noindent\textbf{Scope and priority.}
The exact recurrence and table in Bayoumi--El-Zahar--Khamis (1990),
classical two-dimensional-poset asymptotics and sampling precedents,
and the simple-permutation asymptotic machinery are prior art.
This report presents an all-orders refinement; it makes no first-ever
claim for the leading equivalent. The final 2026 intersection-graph
article and two older subscription articles were identified but their
full final texts were not inspected. All truncation orders are fixed;
no effective onset or growing-order remainder estimate is claimed.
\newpage
\tableofcontents
\newpage

\section{Results and the proof mechanism}\label{sec:results}

A permutation $\pi\in\mathcal S_n$ determines an inversion graph
on $\{1,\ldots,n\}$: vertices $i<j$ are adjacent exactly when
$\pi(i)>\pi(j)$. A \emph{permutation graph} is an abstract graph
isomorphic to such a graph. Throughout, graphs are simple and
unlabeled; a graph and its complement are counted separately unless
they happen to be isomorphic. Write $a_n$ for the number of
isomorphism classes on $n$ vertices, and $r_n$ for the number of
classes with one distinguished vertex, with isomorphisms required
to preserve that vertex. Thus $r_n$ counts vertex orbits across
unrooted graph classes, not $na_n$.

Introduce ordinary formal power series
\begin{equation}\label{eq:series}
 A(x)=\sum_{n\geq1}a_nx^n,\qquad
 C(x)=\sum_{d\geq0}r_{d+1}x^d,\qquad
 F(x)=\sum_{n\geq1}n!x^n.
\end{equation}
Let $T=F^{\langle-1\rangle}$ denote formal compositional inversion,
and put $Q=T\circ A$. Neither convergence nor a positive radius of
convergence is intended for these series.

\begin{theorem}[All fixed orders]\label{thm:main}
Define
\begin{equation}\label{eq:H}
 H(x)=C(x)\frac{xQ'(x)}{Q(x)A'(x)}
       \exp\left(\frac1x-\frac1{Q(x)}\right)
      =\sum_{k\geq0}h_kx^k.
\end{equation}
This is a well-defined series in $\QQ[[x]]$, with $H(0)=1$.
For each fixed integer $R\geq1$, as $n\longrightarrow\infty$,
\begin{equation}\label{eq:main}
 a_n=\frac14\sum_{k=0}^{R-1}h_k(n-k)!
       +O_R((n-R)!).
\end{equation}
The coefficient $h_k$ depends only on $a_1,\ldots,a_{k+2}$ and
$r_1,\ldots,r_{k+1}$. These are finite, exactly enumerable inputs.
Moreover, $k!h_k\in\ZZ$ for every $k\geq0$.
\end{theorem}

The first eight $h_k$ are displayed in Section~\ref{sec:coefficients}.
In particular,
\begin{equation}\label{eq:three}
 a_n=\frac{n!}{4}\left(1-\frac4n-\frac1{n^2}
                           -\frac{97}{3n^3}+O(n^{-4})\right).
\end{equation}
The inverse-factorial and inverse-power coefficient lists are
\emph{different}: the $-94/3$ coefficient multiplies $(n-3)!$ in
\eqref{eq:main}, whereas $-97/3$ multiplies $n^{-3}$ after division
by $n!$.

The leading equivalent gives a natural comparison of sampling laws.
Let $\mathsf P_n$ be the law of the unlabeled inversion graph of a
uniform permutation in $\mathcal S_n$, and let $\mathsf U_n$ be
uniform on the $a_n$ graph classes. With
$d_{\TV}(\mu,\nu)=\sup_E|\mu(E)-\nu(E)|$, we prove
\begin{equation}\label{eq:tvintro}
 d_{\TV}(\mathsf P_n,\mathsf U_n)=\frac4n-\frac{15}{n^2}
 +\frac{169}{3n^3}+\frac{215}{2n^4}+O(n^{-5}).
\end{equation}
This controls events and uniformly bounded observables. It does not,
by itself, transfer expectations of unbounded statistics or relative
errors for rare events.

For $N(X)=\min\{n\geq1:a_n\geq X\}$, the report gives arbitrary
fixed-order smooth approximations and a shrinking rounding window.
In particular, with $L=\log X$, $u=L/W_0(L/e)$, $D=\log u$, and
$c=\frac12\log(2\pi)-\log4$, the explicit center is
\begin{equation}\label{eq:inverseintro}
 v_*(X)=u-\frac12-\frac{c}{D}
         +\frac{97/24-c^2/(2D^2)}{uD}.
\end{equation}
For some constant $K>0$ and all sufficiently large $X$,
\begin{equation}\label{eq:inverseintrobracket}
 \floor{v_*(X)-\frac{K}{u^2D}}+1
 \ \leq N(X)\leq\
 \ceil{v_*(X)+\frac{K}{u^2D}}.
\end{equation}
The constants and onset are not made effective here, and the floor
and ceiling cannot be replaced uniformly by a single unqualified
ceiling of a truncated expression.

\paragraph{Why the proof is not circular.}
The series $A$ occurs on the right of \eqref{eq:H}, but only as a
formal device for collecting stabilized coefficients. For each
fixed $R$ we first prove an approximation involving the known
simple-permutation series and two finite polynomials. Only those
series are submitted to the analytic asymptotic-transfer theorem.
The resulting coefficients stabilize as the finite cutoff grows.
No closure theorem is applied to $A$ before its factorial expansion
has been established.

\section{Structural conventions and established inputs}\label{sec:inputs}

\subsection{Intervals and modular decomposition}
An interval of a permutation is a consecutive set of positions whose
values also form a consecutive set. Its length may be one or the
whole permutation; a \emph{nontrivial proper interval} has length
between $2$ and $n-1$. A permutation of order at least four is
\emph{simple} when it has no nontrivial proper interval.

A module in a graph $G$ is a nonempty vertex set $M$ such that every
vertex outside $M$ is adjacent to all of $M$ or to none of $M$.
Two sets \emph{overlap} when their intersection and both differences
are nonempty. A module is \emph{strong} when it overlaps no module.
The strong modules form the canonical modular decomposition tree.
At its internal nodes, maximal proper strong modules are the child
blocks. The quotient is either edgeless (parallel), complete
(series), or modular-prime. We use \emph{prime} only in this modular
sense, not in the sense of graph split decomposition.

Inflating a permutation replaces each of its entries by an interval
carrying a smaller permutation. Its canonical substitution tree has
simple nodes of arity at least four and increasing or decreasing
nodes of arity at least two. Consecutive monotone nodes with the
same sign are flattened. The standard decomposition and the
simple-permutation generating equation are due to
Albert--Atkinson--Klazar~\cite{AAK}.

The precise bridge from permutations to graphs is important.
Habib--Paul~\cite[Lemma 20, printed p.~55]{HP} identify the strong
modules of an inversion graph with the strong common intervals of
its two-order representation. Consequently the canonical
substitution tree agrees with the graph's modular decomposition
tree: simple nodes become prime quotients, and monotone nodes become
series or parallel nodes. This does \emph{not} say that every graph
module is an interval of every realizer. It says that strong modules
are exactly the canonical interval blocks relevant here.

Permutation graphs are closed under induced subgraphs and graph
substitution. The latter follows directly by inflating a permutation
realizer of the quotient by realizers of the child graphs. In
particular, every graph counted by a context-and-decoration
construction below is still a permutation graph.

\subsection{Four orientation pairs and weighted prime counts}\label{sec:four}
A permutation graph is a comparability graph whose complement is
also a comparability graph. For a prime permutation graph of order
at least four, the graph has two transitive orientations, reverse
to one another, and its complement likewise has two. Each of the
four pairs determines two total orders and hence a normalized
permutation realizer. These classical facts and their use in
permutation-graph enumeration are reviewed in
\cite{BEK,BBFGP}.

Here a \emph{normalized realizer} is an actual permutation in
$\mathcal S_m$ whose inversion graph belongs to the specified
unlabeled graph class. No external labeling or ordering of the
abstract graph's vertices is counted as additional data.

\begin{lemma}[Weighted four-realizer identity]\label{lem:weighted}
If $P$ is a prime permutation graph on $m\geq4$ vertices, then
\[
 f(P):=\#\{\pi\in\mathcal S_m:\Inv(\pi)\simeq P\}
       =\frac4{|\Aut(P)|}.
\]
If $s_m$ is the number of simple permutations of order $m$, then
\begin{equation}\label{eq:weighted}
 \sum_{[P]\text{ prime},\ |P|=m}\frac1{|\Aut(P)|}=\frac{s_m}{4}.
\end{equation}
\end{lemma}
\begin{proof}
Fix a vertex set for $P$. An automorphism preserving both chosen
transitive orientations preserves their union, a total order on
that vertex set, so it is the identity. The automorphism action on
the four orientation pairs is therefore free. Two pairs give the
same normalized permutation exactly when they differ by an
automorphism. This proves the first formula. Simplicity of a
permutation is equivalent to modular primality of its inversion
graph at these orders, by the canonical decomposition. Summing
$f(P)$ gives $s_m$, proving \eqref{eq:weighted}.
\end{proof}

The four realizer operations form a Klein group generated by
inversion $\pi\mapsto\pi^{-1}$ and reverse-complement
$\pi\mapsto\rho\pi\rho$, where $\rho(i)=m+1-i$. They preserve the
unlabeled inversion graph. Complementing the graph is not an
additional equivalence and supplies no extra factor of two.
Orders one, two, and three do not enter the prime series.

\subsection{The factorial asymptotic transfer used here}\label{sec:borinsky}
Let
\begin{equation}\label{eq:S}
 S(z)=\sum_{m\geq4}s_mz^m=\frac{z-z^2}{1+z}-T(z).
\end{equation}
For the factorial class with parameters $\alpha=\beta=1$, write
$\mathcal D f(x)=\sum_{k\geq0}d_kx^k$ when
\[
 [x^n]f(x)=\sum_{k=0}^{R-1}d_k(n-k)!+O_R((n-R)!)
 \quad\text{for every fixed }R.
\]
The basis is $(n-k)!=\Gamma(n+1-k)$, rather than
$\Gamma(n-k)$. Polynomials have asymptotic derivative zero.

Borinsky's composition and inversion theorems
\cite[Theorem 35 and equations (62)--(67)]{Borinsky} imply
\begin{equation}\label{eq:DS}
 \mathcal D S(z)=T'(z)\frac{z}{T(z)}
             \exp\left(\frac1z-\frac1{T(z)}\right).
\end{equation}
In particular $\mathcal D S(0)=e^{-2}$. For a polynomial $B$ with
$B(0)=0$ and $B'(0)=1$, the specialization needed here is
\begin{equation}\label{eq:composition}
 \mathcal D(S\circ B)(x)=\frac{x}{B(x)}
  \exp\left(\frac1x-\frac1{B(x)}\right)(\mathcal D S)(B(x)).
\end{equation}
For a polynomial $C_0$, multiplication gives
\begin{equation}\label{eq:polymult}
 \mathcal D\bigl(C_0(S\circ B)\bigr)=C_0\mathcal D(S\circ B).
\end{equation}
The additional chain-rule term in the general theorem vanishes
because $\mathcal D B=0$. Equations \eqref{eq:DS}--\eqref{eq:polymult}
are established asymptotic results, not identities inferred solely
from formal power-series manipulation. We state precisely these
external inputs and prove the graph-specific finite approximation
next.

\section{Excluding middle intervals}\label{sec:intervals}

\begin{lemma}[Interval tail]\label{lem:intervals}
For every fixed integer $R\geq1$,
\begin{equation}\label{eq:intervaltail}
 \sum_{k=R+2}^{n-R}\frac{(n-k+1)^2}{\binom nk}=O_R(n^{-R}).
\end{equation}
Consequently at most $O_R((n-R)!)$ permutations of order $n$
have an interval of some length in $[R+2,n-R]$.
\end{lemma}
\begin{proof}
For each interval of $k$ positions, its image set is uniform over
all $\binom nk$ subsets of size $k$. Exactly $n-k+1$ of these
subsets are intervals of values. There are $n-k+1$ intervals of
positions, so the expected number of length-$k$ intervals is
\[
 E_k=\frac{(n-k+1)^2}{\binom nk}.
\]
For $k\leq n/4$,
\[
 \frac{E_{k+1}}{E_k}
 =\frac{(k+1)(n-k)}{(n-k+1)^2}
 \leq\frac{k+1}{n-k+1}\leq\frac12
\]
for all sufficiently large $n$. The lower-end tail is therefore
bounded by a constant times $E_{R+2}=O_R(n^{-R})$.
At the upper end put $\ell=n-k$ and
$b_\ell=(\ell+1)^2/\binom n\ell$. Then
\[
 \frac{b_{\ell+1}}{b_\ell}
 =\frac{(\ell+2)^2}{(\ell+1)(n-\ell)}.
\]
For $1\leq\ell\leq n/4$ this is at most $1/2$ once $n$ is large.
The upper-end tail starts at $b_R=O_R(n^{-R})$.
Throughout the remaining central band,
$\binom nk\geq\binom n{\lfloor n/4\rfloor}$, up to an immaterial
shift of one in a boundary index. This grows exponentially in $n$,
whereas the numerator and number of terms are polynomial. The
central sum is exponentially small. This proves
\eqref{eq:intervaltail}. A union bound and multiplication by $n!$
complete the proof, since $n!n^{-R}=O_R((n-R)!)$.
\end{proof}

\begin{definition}
A permutation is $R$-good if it has no interval of a length in
$[R+2,n-R]$. All goodness statements below have $R$ fixed and
$n$ sufficiently large in terms of $R$.
\end{definition}

The argument is a probability estimate on uniform \emph{permutations}.
It is not an assertion that such permutations induce uniform graph
classes. Later, an exceptional graph class will be charged to one
of its exceptional realizers; this elementary surjection bound is
all that is needed to transfer the error count to graphs.

\section{Recovering one large prime node}\label{sec:core}

\begin{lemma}[Large simple node]\label{lem:large}
The canonical substitution tree of an $R$-good permutation has a
unique internal node of arity greater than $R+1$. That node is
simple. Its subtree has at least $n-R+1$ leaves, each of its children
has at most $R+1$ leaves, and at most $R-1$ leaves lie outside it.
\end{lemma}
\begin{proof}
Every proper canonical subtree is an interval, so its leaf count
is at most $R+1$ or at least $n-R+1$. Starting at the root, follow
a child of the latter size for as long as one exists. Such a child
is unique, because two would have total size exceeding $n$. At the
terminal node $v$, the subtree has $N\geq n-R+1$ leaves and all
children have at most $R+1$ leaves.

The node cannot be monotone. Otherwise any consecutive union of its
children is also a permutation interval. Accumulate children from
one end until their total size first exceeds $R+1$. The resulting
interval has size between $R+2$ and $2R+2$. For $n>3R+2$, this lies
in the forbidden band and is proper. Notice that the union need not
be a canonical subtree: being a proper \emph{permutation interval}
is the property used in the contradiction.

Thus $v$ is simple and its arity is at least $N/(R+1)$, eventually
greater than $R+1$. Any descendant internal node lies inside a child
with at most $R+1$ leaves. There are at most $R-1$ leaves outside
$v$, so every ancestor has at most $R$ children: one on the path to
$v$ and at most $R-1$ off that path. Every off-path node is likewise
bounded. This proves uniqueness and all size assertions.
\end{proof}

To turn this into an unlabeled product, the large inserted object
must be intrinsically recoverable in every bounded context.

\begin{lemma}[Connected and coconnected insertion]\label{lem:insertion}
Suppose a graph $H$ is substituted at one vertex of a graph $K$.
If $H$ and its complement are both connected, the set $M$ of
inserted vertices is a strong module of the resulting graph.
\end{lemma}
\begin{proof}
It is a module by construction. Suppose a module $N$ overlaps $M$.
Let $X=M\cap N$ and choose $y\in N\setminus M$. Since $M$ is a
module, $y$ has constant adjacency $\eta\in\{0,1\}$ to $M$.
For any $z\in M\setminus N$, the module property of $N$ forces
$z$ to have the same adjacency to every $x\in X$ as it has to $y$.
Thus every edge across the nontrivial cut $(X,M\setminus N)$ has
value $\eta$. If $\eta=0$, $H$ is disconnected; if $\eta=1$, its
complement is disconnected. Both contradict the hypothesis.
\end{proof}

A prime quotient of order at least four is connected and
coconnected. Inflating all its vertices by nonempty graphs preserves
both properties: paths in the quotient connect different blocks,
and a vertex in another adjacent block connects vertices within a
block if necessary. Apply the same argument in the complement.
Thus Lemma~\ref{lem:insertion} applies to every prime-rooted
inflation used here, regardless of the surrounding context.

\begin{proposition}[Unlabeled context recovery]\label{prop:context}
Fix $R$. Consider the following construction at order $n$:
\begin{enumerate}
\item choose a vertex-rooted permutation graph on $d+1$ vertices,
      where $0\leq d\leq R-1$;
\item choose a prime permutation-graph quotient of order $m>R+1$;
\item decorate each quotient vertex by a nonempty permutation graph
      of order at most $R+1$, modulo quotient automorphisms;
\item insert the decorated quotient at the marked context vertex.
\end{enumerate}
For sufficiently large $n$, the output graph recovers exactly these
data, up to the stated isomorphisms. Every graph realized by an
$R$-good permutation occurs among the outputs.
\end{proposition}
\begin{proof}
The inserted vertex set is a strong module by
Lemma~\ref{lem:insertion}. Within it, the maximal proper strong
modules are the decorated ports and the quotient is the chosen
prime graph. All decoration nodes have arity at most $R+1$, as do
all nodes of the context, whose order is at most $R$. Hence the
chosen quotient is the unique node of arity greater than $R+1$
in the output's canonical modular tree. It is intrinsically
recognizable, including when the context consists only of its root.

An isomorphism between output graphs must preserve this unique node
and its strong module. The induced graph on that module recovers
the inner decorated quotient, with ports identified up to quotient
automorphism. Collapsing it to one marked vertex recovers the
context up to root-preserving isomorphism. Conversely each listed
isomorphism of data gives an isomorphism of outputs. There is no
extra division by a context root stabilizer and no choice of a
transitive orientation at the context root.

Finally Lemma~\ref{lem:large} and the Habib--Paul correspondence
identify exactly such data inside every $R$-good realizer. Taking
induced subgraphs and collapsing interval blocks preserves the
permutation-graph property. Thus its graph lies among the outputs.
\end{proof}

Let $\mathcal U_R(n)$ be this exact subclass of graph classes. It
may contain graphs with some bad realizers. No converse goodness
assertion is needed. The two inclusions used are simply
\[
 \mathcal U_R(n)\subseteq\{\text{all graph classes of order }n\},
 \qquad
 \Inv(\{R\text{-good permutations}\})\subseteq\mathcal U_R(n).
\]
Every class outside $\mathcal U_R(n)$ therefore has only bad
realizers. Selecting one realizer per such class and applying
Lemma~\ref{lem:intervals} gives
\begin{equation}\label{eq:subclasserror}
 0\leq a_n-|\mathcal U_R(n)|=O_R((n-R)!).
\end{equation}
The rooted context factor is precisely $r_{d+1}$. Port decorations,
by contrast, are \emph{unrooted}, since outside adjacency is uniform
throughout a port module.

\section{Burnside weights and negligible symmetric quotients}\label{sec:symmetry}

Define finite polynomials
\begin{equation}\label{eq:BC}
 B_R(x)=\sum_{j=1}^{R+1}a_jx^j,\qquad
 C_R(x)=\sum_{d=0}^{R-1}r_{d+1}x^d.
\end{equation}
For a fixed prime graph $P$ with $m$ vertices, graph decorations of
its ports are acted on by $\Aut(P)$. Burnside's lemma says that
their orbit generating series is the average of the fixed-decoration
series. The identity contribution is
$B_R(x)^m/|\Aut(P)|$. Summation over prime classes using
\eqref{eq:weighted} gives exactly $S(B_R(x))/4$ for these identity
contributions.

The identity contribution is a rational weight, not in general an
integer count of a graph family. We must bound the missing
nonidentity Burnside contributions before using it in an
asymptotic formula.

\begin{lemma}[Symmetry bound]\label{lem:symmetry}
The number of permutations of order $m$ fixed by at least one
nonidentity Klein-group operation is at most
\begin{equation}\label{eq:J}
 J_m=\exp\left(\frac m2\log m+O(m)\right).
\end{equation}
For fixed $R$, all decorated symmetric prime-quotient contributions
at total graph order $n$, including bounded contexts, have total
size at most
\begin{equation}\label{eq:symdecor}
 \exp\left(\frac n2\log n+O_R(n)\right)
       =o((n-R)!).
\end{equation}
\end{lemma}
\begin{proof}
Inversion-fixed permutations are involutions. The product of
inversion and reverse-complement is fixed when
$\pi=\rho\pi^{-1}\rho$, equivalently $(\rho\pi)^2=1$, so it has
the same count. Reverse-complement-fixed permutations commute with
$\rho$. There are exactly
$2^{\lfloor m/2\rfloor}\lfloor m/2\rfloor!$ of them: permute the
reversal pairs and independently choose their directions, fixing
the middle point at odd order.

If $I_m$ counts involutions, its exponential generating function is
$\exp(z+z^2/2)$. Positivity of its coefficients gives, at
$t=\sqrt m$,
\[
 I_m\leq m!\,\frac{\exp(t+t^2/2)}{t^m}
      =\exp\left(\frac m2\log m+O(m)\right).
\]
The reverse-complement count satisfies the same bound by Stirling's
formula. A union bound proves \eqref{eq:J}.

If $P$ has a nontrivial automorphism, Lemma~\ref{lem:weighted}
shows that it has fewer than four distinct realizers. Its Klein
orbit therefore has a nontrivial stabilizer. Such prime classes
can be injected into the union just counted by choosing one of
their realizers.

Put $b_R=\sum_{j=1}^{R+1}a_j$. The number of decoration assignments
at $m$ distinguishable ports, ignoring total size, is at most
$b_R^m$. This also bounds each fixed-assignment term, and
$|\Aut(P)|\leq4$. Summing over $m\leq n$ and multiplying by the
bounded number $\sum_{d<R}r_{d+1}$ of contexts gives the first
bound in \eqref{eq:symdecor}. Since
$\log((n-R)!)=n\log n-n+O_R(\log n)$, the second assertion follows.
\end{proof}

\begin{theorem}[Finite-polynomial approximation]\label{thm:finite}
For every fixed $R\geq1$,
\begin{equation}\label{eq:finite}
 a_n=\frac14[x^n]\,C_R(x)S(B_R(x))+O_R((n-R)!).
\end{equation}
\end{theorem}
\begin{proof}
Proposition~\ref{prop:context} makes the class construction an
ordinary unlabeled product of contexts and decorated quotients.
Burnside's identity terms give the displayed coefficient.
Lemma~\ref{lem:symmetry} bounds all nonidentity terms by an error
smaller than every fixed factorial order. The series $S$ also
contains quotients with arity $m\leq R+1$; these contribute only
orders at most $(R+1)^2+R-1$ after decoration and context insertion,
so they disappear at sufficiently large $n$. Finally use
\eqref{eq:subclasserror}.
\end{proof}

The theorem has no hidden hypothesis about the unknown growth of
$a_n$ or $r_n$: for fixed $R$, they enter only as finitely many
integer coefficients of $B_R$ and $C_R$. The constants may depend
on those finite inputs and on $R$.

\section{Factorial transfer and all orders}\label{sec:transfer}

We now prove Theorem~\ref{thm:main}. Apply
\eqref{eq:composition} and \eqref{eq:polymult} to
Theorem~\ref{thm:finite}. For a fixed $R$, the coefficients through
degree $R-1$ are generated by
\begin{equation}\label{eq:HR}
 H_R(x)=C_R(x)\frac{x}{B_R(x)}
       \exp\left(\frac1x-\frac1{B_R(x)}\right)
       (\mathcal D S)(B_R(x)).
\end{equation}
Precisely,
\[
 a_n=\frac14\sum_{k=0}^{R-1}[x^k]H_R(x)(n-k)!
             +O_R((n-R)!).
\]
The finite graph approximation and the factorial-transfer remainder
have the same order and can be combined.

Because $A-B_R=O(x^{R+2})$ and both have leading coefficient one,
\[
 \frac1A-\frac1{B_R}=O(x^R),\qquad
 \frac{x}{A}-\frac{x}{B_R}=O(x^{R+1}).
\]
Also $C-C_R=O(x^R)$, and ordinary formal substitution into
$\mathcal D S$ changes no coefficient below degree $R$. Hence
\begin{equation}\label{eq:stabilize}
 H_R(x)\equiv C(x)\frac{x}{A(x)}
       \exp\left(\frac1x-\frac1{A(x)}\right)
       (\mathcal D S)(A(x))\pmod{x^R}.
\end{equation}
Substitute \eqref{eq:DS}. The exponentials combine by cancellation
of $1/A$, and $T'(A)=Q'/A'$ by the formal chain rule. Thus the
right side of \eqref{eq:stabilize} is exactly \eqref{eq:H}.
This proves \eqref{eq:main} at every fixed order.

For well-definedness, $a_1=1$ and $a_2=2$, while
$F(x)=x+2x^2+\cdots$. Solving $F(Q)=A$ gives
$Q=x+O(x^3)$. Therefore
\[
 E(x):=\frac1x-\frac1{Q(x)}\in x\QQ[[x]],
 \qquad \frac{xQ'}{QA'}\in1+x\QQ[[x]],
\]
and $C(0)=r_1=1$. In particular the constant terms in the
simple-permutation transfer, $e^{-2}$ from $\mathcal D S$ and
$e^2$ from $B_R=x+2x^2+\cdots$, cancel. This explains why the
leading graph constant is $1/4$ rather than $e^{-2}/4$.

The degree-$k$ exponent coefficient uses $Q$ only through degree
$k+2$. Since $F(Q)=A$ has a unit linear term, these coefficients
use only $a_1,\ldots,a_{k+2}$. The prefactor needs no more data,
and its $C$ factor uses $r_1,\ldots,r_{k+1}$. This proves the
finite dependency assertion.

\begin{corollary}[Denominators]\label{cor:denominator}
Every coefficient in \eqref{eq:H} satisfies $k!h_k\in\ZZ$.
\end{corollary}
\begin{proof}
The inverse of a series $x+O(x^2)$ with integer coefficients has
integer coefficients: solving successive coefficients involves no
division other than by its linear coefficient one. Thus $T$, $Q$,
and $Q/x$ have integer coefficients. The units $Q/x$ and $A'$ have
constant term one, so
\[
 P(x):=C(x)\frac{xQ'(x)}{Q(x)A'(x)}\in\ZZ[[x]].
\]
Moreover $Q=x+O(x^3)$ implies $E\in x\ZZ[[x]]$. In degree $k$,
$\exp E=\sum_{j=0}^kE^j/j!$ has denominator dividing $k!$.
Multiplication by $P$ preserves this property, because $j!$ divides
$k!$ for $j\leq k$.
\end{proof}

\begin{remark}[The meaning of all orders]\label{rem:fixed}
The quantifiers are: for each fixed $R$ there exist a constant and
an onset depending on $R$. The proof does not bound their growth
with $R$ and does not justify substituting $R=R(n)$, choosing a
least term, claiming convergence or Borel summability of $H$, identifying all
exponentially small sectors, or estimating an
exponentially small remainder by optimally truncating the series.
Finite exact coefficient algorithms and an all-orders asymptotic
statement do not establish any of those stronger conclusions.
\end{remark}

\section{Exact coefficients and finite certification}\label{sec:coefficients}

The necessary small graph data are as follows. The unrooted and rooted
columns were obtained by two separately implemented exhaustive
canonicalization routes; no source-table value beyond these orders
is used in the asymptotic calculation.
\begin{center}
\begin{tabular}{r r r r r}
\toprule
$n$ & $a_n$ & $r_n$ & $s_n$ & Prime graph classes\\
\midrule
1&1&1&0&0\\
2&2&2&0&0\\
3&4&6&0&0\\
4&11&20&2&1\\
5&33&89&6&3\\
6&142&504&46&16\\
7&776&3776&338&92\\
8&5699&34620&2926&772\\
9&50723&368683&28146&7122\\
\bottomrule
\end{tabular}
\end{center}
The zero values in the prime columns below order four reflect the
convention that the prime series starts at four. Both enumerations
check the unrooted and vertex-rooted columns; the independent
refined-color route additionally supplies the displayed simple and
prime counts. In particular all inputs for $h_0,\ldots,h_7$,
namely $a_1,\ldots,a_9$ and $r_1,\ldots,r_8$, have two exhaustive
replays. The extra $r_9$ does not by itself determine $h_8$, which
would also require $a_{10}$.

Solving $F(Q)=A$ gives
\begin{align*}
 Q(x)={}&x-2x^3-5x^4-31x^5-180x^6-1418x^7\\
       &{}-12327x^8-121659x^9+O(x^{10}).
\end{align*}
The coefficient table is
\begin{center}
\begin{tabular}{r r r}
\toprule
$k$ & $h_k$ & $k!h_k$\\
\midrule
0&$1$&$1$\\
1&$-4$&$-4$\\
2&$-1$&$-2$\\
3&$-94/3$&$-188$\\
4&$-769/6$&$-3076$\\
5&$-19969/15$&$-159752$\\
6&$-531812/45$&$-8508992$\\
7&$-40114096/315$&$-641825536$\\
\bottomrule
\end{tabular}
\end{center}

For clarity about the first correction, one can calculate directly:
\[
 Q=x-2x^3+O(x^4),\quad E=-2x+O(x^2),\quad
 C\frac{xQ'}{QA'}=1-2x+O(x^2).
\]
Hence $H=1-4x+O(x^2)$. Both the rooted context and the prime
inflation contribute; replacing $C$ by an unrooted series would
produce incorrect higher coefficients.

\subsection{Two exact formal calculation routes}
One route computes $Q$ recursively from $F(Q)=A$ and evaluates
\eqref{eq:H} with rational truncated-series arithmetic. A second
route obtains $T$ by Lagrange inversion,
\begin{equation}\label{eq:lagrange}
 [z^m]T(z)=\frac1m[w^{m-1}]\left(\frac{w}{F(w)}\right)^m,
\end{equation}
then evaluates \eqref{eq:HR} for successive finite cutoffs.
The factor $e^{-2}$ in \eqref{eq:DS} can be carried symbolically
and canceled with the constant $e^2$ in the composition factor.
Every remaining calculation is rational. The routes agree through
$h_7$, and the finite-cutoff route checks the stabilization in
\eqref{eq:stabilize} for $R=1,\ldots,8$.

\subsection{Two graph enumeration routes}
Each route enumerates all $n!$ permutations, constructs its inversion
graph, and computes an exact isomorphism key. The first partitions
vertices by degree (and, for rooted keys, root status), then minimizes
an adjacency encoding over all permutations within the cells.
For each unrooted representative it counts distinct rooted keys.
The second refines invariant color classes before exhaustively
minimizing within them; all minimizing labelings recover the
vertex orbits and automorphism count. Both are exhaustive rather
than randomized isomorphism heuristics.

Additional finite checks verify the weighted identity
$|\Aut(P)|f(P)=4$ and equality of the realizer set with the Klein
orbit for every prime graph through order nine. Direct subset-module
checks at smaller orders confirm the prime/simple correspondence.
These computations are useful normalization tests. They do not
replace the structural theorems or certify an asymptotic remainder
constant.

\section{Threshold inversion with rounding preserved}\label{sec:inverse}

\subsection{Strict increase and ordinary inverse powers}
Adding an isolated vertex gives an injection from graph classes of
order $n$ to those of order $n+1$: remove one isolated vertex to
recover the original class. The complete graph on $n+1$ vertices
is outside the image for $n\geq1$. Thus
\begin{equation}\label{eq:monotonicity}
 a_{n+1}>a_n\qquad(n\geq1).
\end{equation}
The sequence tends to infinity by Theorem~\ref{thm:main}, so $N(X)$
is well defined for all real $X>0$.

Let $(n)_k=n(n-1)\cdots(n-k+1)$, with $(n)_0=1$, and define a
formal inverse-power series
\begin{equation}\label{eq:U}
 U(y)=\sum_{k\geq0}h_k\frac{y^k}{\prod_{j=0}^{k-1}(1-jy)}.
\end{equation}
At each degree this is a finite sum. Its beginning is
\begin{equation}\label{eq:Ubegin}
 U(y)=1-4y-y^2-\frac{97}{3}y^3-\frac{1339}{6}y^4
                         -\frac{11603}{5}y^5+O(y^6).
\end{equation}
Using a formal logarithm and Stirling's expansion gives, in the
usual fixed-order asymptotic sense,
\begin{align}\label{eq:logexp}
 \log a_n={}&n(\log n-1)+\frac12\log n+c
            -\frac{47}{12n}-\frac9{n^2}+O(n^{-3}),\\
 &\hspace{15mm}c=\frac12\log(2\pi)-\log4.\nonumber
\end{align}
All further inverse-power coefficients follow from \eqref{eq:U},
$\log U$, and the Bernoulli-number terms in Stirling's formula.
For example $-47/12=1/12-4$ and $-9=-1-(-4)^2/2$.

\subsection{An arbitrary fixed-order smooth approximation}
For a fixed $R\geq1$, set
\begin{equation}\label{eq:AR}
 \mathcal A_R(v)=\frac{\Gamma(v+1)}4
       \sum_{k=0}^{R-1}\frac{h_k}{(v)_k},
 \qquad (v)_k=\prod_{j=0}^{k-1}(v-j).
\end{equation}
For sufficiently large real $v$, the finite factor is positive and
\begin{equation}\label{eq:slope}
 (\log\mathcal A_R)'(v)=\log v+O_R(v^{-1})>0.
\end{equation}
For all sufficiently large $X$, let $\nu_R(X)$ be the unique large
solution of $\mathcal A_R(v)=X$. The qualifier ``large'' excludes
irrelevant behavior of the rational factor near its poles at small
integers. The factorial expansion implies
\begin{equation}\label{eq:logerror}
 \log a_n-\log\mathcal A_R(n)=O_R(n^{-R}).
\end{equation}

\begin{theorem}[Threshold window]\label{thm:threshold}
For each fixed $R\geq1$, there is a positive function
\[
 \eps_R(X)=O_R\left(\frac{\nu_R(X)^{-R}}{\log\nu_R(X)}\right)
\]
such that, for all sufficiently large $X$,
\begin{equation}\label{eq:window}
 \floor{\nu_R(X)-\eps_R(X)}+1
       \leq N(X)\leq\ceil{\nu_R(X)+\eps_R(X)}.
\end{equation}
In particular $N(X)=\ceil{\nu_R(X)}$ whenever $\nu_R(X)$ stays
farther than $\eps_R(X)$ from every integer.
\end{theorem}
\begin{proof}
Choose $K_R$ and $n_R$ so that the absolute error in
\eqref{eq:logerror} is at most $K_Rn^{-R}$ for integers $n\geq n_R$.
Put $v=\nu_R(X)$. On $[v/2,2v]$, equation~\eqref{eq:slope} gives
$(\log\mathcal A_R)'\geq\frac12\log v$ once $v$ is large.
Choose
\[
 \eps_R=D_R\frac{v^{-R}}{\log v},\qquad D_R>2^{R+1}K_R.
\]
If an integer $n$ in that interval satisfies $n\leq v-\eps_R$,
then
\[
 \log\mathcal A_R(n)-\log X
 \leq-\tfrac12\eps_R\log v
 <-2^RK_Rv^{-R}\leq-K_Rn^{-R}.
\]
Together with \eqref{eq:logerror}, this yields $a_n<X$.
The same argument with the opposite sign gives $a_n>X$ for
integers $n\geq v+\eps_R$ in the interval. Monotonicity
\eqref{eq:monotonicity} extends these comparisons to integers
outside it, using nearby integers still inside the interval.
Thus every integer at or below the lower edge is too small, and
every integer at or above the upper edge is sufficient. Taking
floor and ceiling yields \eqref{eq:window}, including integer-edge
cases because the comparison inequalities are strict.
\end{proof}

\subsection{The explicit Lambert W center}
Let $W_0$ be the principal real Lambert function, so
$W_0(z)e^{W_0(z)}=z$. With
\begin{equation}\label{eq:u}
 L=\log X,\qquad u=\frac{L}{W_0(L/e)},\qquad D=\log u,
\end{equation}
one has $u(\log u-1)=L$. Set
\[
 g(v)=v(\log v-1)+\tfrac12\log v+c-\frac{47}{12v}.
\]
Equation~\eqref{eq:logexp} says $\log a_n=g(n)+O(n^{-2})$.
Write $\delta_0=-\frac12-c/D$. The terms of order one in
$g(u+\delta_0)-L$ cancel because
$D\delta_0+D/2+c=0$. Taylor expansion gives the remaining first
inverse-power residual
\begin{equation}\label{eq:residual}
 g(u+\delta_0)-L
 =\frac{\delta_0^2/2+\delta_0/2-47/12}{u}+O(u^{-2})
 =\frac{-97/24+c^2/(2D^2)}u+O(u^{-2}).
\end{equation}
Canceling this with a second shift divided by the leading slope
$D$ gives the $v_*$ in \eqref{eq:inverseintro}. Since
$g'(v)=\log v+O(v^{-1})$, its residual is $O(u^{-2})$ and its
smooth inverse error is $O(u^{-2}/D)$.

The same strict-comparison argument used for
Theorem~\ref{thm:threshold}, now using
$\log a_n-g(n)=O(n^{-2})$, proves
\eqref{eq:inverseintrobracket}. In particular the less refined
center
\[
 u-\frac12-\frac{c}{D}
\]
has inverse error $O(1/(uD))$ before rounding. Under the alternate
normalization $L_4=\log(4X)$ and $t=L_4/W_0(L_4/e)$, its first
shift is often written
\[
 t-\frac{\log(2\pi t)}{2\log t}.
\]
These are compatible after re-expansion; one must not mix $\log X$
and $\log(4X)$ while retaining the same constant $c$.

\subsection{What the inverse does and does not certify}
The coefficients and formal shifts can be computed to any fixed
order. Once $R$ is chosen, a large root of \eqref{eq:AR} can be
approximated numerically, for example by Newton iteration initialized
at a Lambert-$W$ center. This is a method for evaluating the smooth
approximation, not a certified algorithm for the exact integer
threshold at an arbitrary supplied $X$: no explicit $K_R$ or
uniformly effective onset has been obtained here.

A staircase with jumps of one cannot be approximated uniformly by a
continuous expression to $o(1)$. Near an integer center, the two
neighboring integers in the bracket are genuine possibilities.
Away from the shrinking boundary window the ceiling is stable.
Neither additional formal coefficients nor higher floating-point
precision remove this logical distinction.

\section{All orders comparison of the sampling laws}\label{sec:tv}

Let $\mathcal G_n$ be the finite set of unlabeled permutation graphs
of order $n$, and write $f(G)$ for the number of normalized
permutation realizers of $G$. Then
\begin{equation}\label{eq:laws}
 \mathsf P_n(G)=\frac{f(G)}{n!},\qquad
 \mathsf U_n(G)=\frac1{a_n},\qquad
 \sum_{G\in\mathcal G_n}f(G)=n!.
\end{equation}
Put $M_{4,n}=\#\{G\in\mathcal G_n:f(G)=4\}$. We first derive an
all-orders expansion for this exact fiber-size class, then transfer
it to total variation.

\subsection{Marked context realizers and exact multiplicativity}
For a rooted graph $(L,v)$, define its marked-realizer count by
\begin{equation}\label{eq:t}
 t(L,v)=\#\{(\kappa,i):(\Inv(\kappa),i)\simeq(L,v)\}
       =f(L)|\Aut(L)\cdot v|.
\end{equation}
For each unmarked realizer $\kappa$, the allowed marked positions
are exactly the images of the automorphism orbit of $v$. This
proves the second equality. Root positions are counted separately
even when the underlying permutation is unchanged.

\begin{lemma}[Exact fiber product]\label{lem:fiberproduct}
Consider a graph $G\in\mathcal U_R(n)$ with asymmetric prime
quotient $P$, port graphs $(D_p)_{p\in V(P)}$, and outer rooted
context $(L,v)$. Then
\begin{equation}\label{eq:fiberproduct}
 f(G)=4\,t(L,v)\prod_{p\in V(P)}f(D_p).
\end{equation}
\end{lemma}
\begin{proof}
Let $H=P[D_p:p\in V(P)]$ be the inserted inner graph. Its child
modules are its maximal proper strong modules. In every realizer
of $H$ they are canonical common-interval blocks. Collapsing them
gives one of the four distinct normalized simple realizers of $P$.
For each one the identification with the ports of $P$ is unique,
because $P$ has no automorphisms. Standardizing each interval gives
an independent choice of a realizer of $D_p$. Conversely these
choices reconstruct the permutation by inflation. Collapse and
standardization recover the inputs, so
$f(H)=4\prod_pf(D_p)$.

In $G$, the vertices of $H$ form the distinguished strong module
having the uniquely large prime quotient. This follows from
Lemma~\ref{lem:insertion} even when the context root belongs to a
series or parallel node. Every realizer of $G$ therefore identifies
that module as a strong common interval. Collapsing it, retaining
its position, and standardizing its interior gives a marked context
realizer counted by $t(L,v)$ and an inner realizer counted by
$f(H)$. Conversely their marked inflation reconstructs the
realizer of $G$. The intrinsic recovery proves this is a bijection.
\end{proof}

Every factor after the initial four in \eqref{eq:fiberproduct}
is a positive integer. An asymmetric-core output has exactly four
realizers if and only if every port graph has one realizer and
its context has one marked realizer. The latter condition need not
force a one-vertex context.

Define finite counts
\begin{align}\label{eq:bc}
 b_j&=\#\{G\in\mathcal G_j:f(G)=1\},\\
 c_j&=\#\{(L,v):|L|=j,\ t(L,v)=1\},\nonumber
\end{align}
and formal series
\begin{equation}\label{eq:JK}
 J(x)=\sum_{j\geq1}b_jx^j,\qquad
 K(x)=\sum_{d\geq0}c_{d+1}x^d.
\end{equation}
Let $J_R$ and $K_R$ truncate these series at degrees $R+1$ and
$R-1$, respectively.

\begin{theorem}[Four-realizer factorial expansion]\label{thm:M4}
Set $Q_J=T\circ J$ and
\begin{equation}\label{eq:H4}
 H_4(x)=K(x)\frac{xQ_J'(x)}{Q_J(x)J'(x)}
       \exp\left(\frac1x-\frac1{Q_J(x)}\right)
       =\sum_{k\geq0}g_kx^k.
\end{equation}
For every fixed $R\geq1$,
\begin{equation}\label{eq:M4}
 M_{4,n}=\frac14\sum_{k=0}^{R-1}g_k(n-k)!
                   +O_R((n-R)!).
\end{equation}
The coefficient $g_k$ uses only $b_1,\ldots,b_{k+2}$ and
$c_1,\ldots,c_{k+1}$, and $k!g_k\in\ZZ$.
\end{theorem}
\begin{proof}
The same finite-core bounds as in Theorem~\ref{thm:finite}, combined
with Lemma~\ref{lem:fiberproduct}, give
\begin{equation}\label{eq:M4finite}
 M_{4,n}=\frac14[x^n]K_R(x)S(J_R(x))+O_R((n-R)!).
\end{equation}
Indeed graph classes outside $\mathcal U_R(n)$ number
$O_R((n-R)!)$, hence bound the excluded four-realizer classes too.
Symmetric large quotients contribute only
$\exp((n/2)\log n+O_R(n))$ classes; this bounds both any genuine
four-realizer contributions from them and their identity terms in
the displayed approximation. At asymmetric quotients there are
exactly four simple realizers and no port identifications; the
fiber product selects precisely the port and context counts in
$J_R$ and $K_R$. Bounded quotients disappear at large $n$ as before.

Apply the same polynomial composition theorem to \eqref{eq:M4finite}
and stabilize the coefficients, with $J,K$ replacing $A,C$.
Since $b_1=1$, $b_2=2$, and $c_1=1$, one has
$Q_J=x+O(x^3)$ and $H_4(0)=1$. The finite-dependency and denominator
arguments of Section~\ref{sec:transfer} apply verbatim. No
asymptotic property of $J$ or $K$ was assumed before finite transfer.
\end{proof}

\subsection{The second coefficient and exact finite inputs}
At order three, the empty and complete graphs have one realizer
each; the edge-plus-isolated-vertex graph and its complement have
two each. Thus $(b_1,b_2,b_3)=(1,2,2)$. Both rooted contexts of
order two have a root orbit of size two, giving $c_1=1,c_2=0$.
At $R=2$, \eqref{eq:M4finite} reduces to
\[
 M_{4,n}=\tfrac14[x^n]S(x+2x^2+2x^3)+O((n-2)!).
\]
In the composition transfer the three factors begin
\begin{align*}
 x/J_2(x)&=1-2x+O(x^2),\\
 \exp(1/x-1/J_2(x))&=e^2(1-2x+O(x^2)),\\
 (\mathcal D S)(J_2(x))&=e^{-2}(1-4x+O(x^2)).
\end{align*}
Hence $g_1=-8$ and
\begin{equation}\label{eq:M4first}
 M_{4,n}=\frac{n!}{4}\left(1-\frac8n+O(n^{-2})\right).
\end{equation}

Full fiber enumeration and explicit vertex-orbit checks give
\begin{center}
\begin{tabular}{r r r r r r r r r}
\toprule
$j$&1&2&3&4&5&6&7&8\\
\midrule
$b_j$&1&2&2&4&2&8&4&22\\
$c_j$&1&0&0&0&0&0&2&0\\
\bottomrule
\end{tabular}
\end{center}
The nonzero $c_7$ is an explicit warning that the outer factor $K$
cannot be discarded in an all-orders formula. The companion checks
these counts and obtains
\[
 (g_0,\ldots,g_6)=
 \left(1,-8,22,-\frac{230}{3},\frac{382}{3},
                 -\frac{22468}{15},-\frac{283418}{45}\right).
\]
The finite-polynomial and direct $Q_J$ routes agree. A separate
finite structural test compared the fiber product against full
permutation fibers on all 3,004 asymmetric $R=2$ constructions
through order eight, finding no discrepancy or construction
collision. That test is additional evidence for the finite
implementation; the all-orders claim rests on the bijection and
error bounds above.

\subsection{The exact sign split and the total variation series}
A class with $f(G)<4$ has a realizer whose Klein orbit has fewer
than four elements, hence a nontrivial stabilizer. Selecting one
such realizer per class and using Lemma~\ref{lem:symmetry} bounds
the number of these classes by
$\exp((n/2)\log n+O(n))$. Their total contribution to either
sampling law is smaller than every inverse power of $n$.

The first correction in \eqref{eq:three} implies
$4<n!/a_n<5$ eventually. Since fiber sizes are integers,
$\mathsf P_n(G)-\mathsf U_n(G)$ is negative exactly for $f(G)\leq4$.
Total variation is the negative mass with its sign reversed.
Write $A_n=4a_n/n!$ and $B_n=4M_{4,n}/n!$. For every fixed $j$,
\begin{equation}\label{eq:TVexact}
 d_{\TV}(\mathsf P_n,\mathsf U_n)
 =M_{4,n}\left(\frac1{a_n}-\frac4{n!}\right)+o(n^{-j})
 =B_n(A_n^{-1}-1)+o(n^{-j}).
\end{equation}
This uses only the integer sign threshold, not divisibility of all
fiber sizes by four. The latter is false: finite fibers of sizes
three, five, and seven occur.

Define the ordinary-power version of the $H_4$ series by
\begin{equation}\label{eq:V}
 V(y)=\sum_{k\geq0}g_k\frac{y^k}{\prod_{j=0}^{k-1}(1-jy)}.
\end{equation}
Then $A_n$ and $B_n$ have the fixed-order inverse-power expansions
$U(1/n)$ and $V(1/n)$, respectively. Since $U(0)=1$, finite series
inversion in \eqref{eq:TVexact} proves the following.

\begin{theorem}[All-orders total variation]\label{thm:tv}
Put
\begin{equation}\label{eq:TVseries}
 Z(y)=V(y)(U(y)^{-1}-1)=\sum_{k\geq1}z_ky^k.
\end{equation}
For every fixed integer $q\geq1$,
\begin{equation}\label{eq:TVall}
 d_{\TV}(\mathsf P_n,\mathsf U_n)
       =\sum_{k=1}^{q}z_kn^{-k}+O_q(n^{-q-1}).
\end{equation}
In particular,
\[
 d_{\TV}(\mathsf P_n,\mathsf U_n)=\frac4n-\frac{15}{n^2}
           +\frac{169}{3n^3}+\frac{215}{2n^4}+O(n^{-5}).
\]
\end{theorem}

To check the second coefficient alone, use
$A_n^{-1}-1=4/n+17/n^2+O(n^{-3})$ and
$B_n=1-8/n+O(n^{-2})$. Their product gives $4/n-15/n^2$.
The uncertainty of order $n^{-2}$ in $B_n$ is multiplied by
$O(n^{-1})$, so it cannot affect this second coefficient.

The complete finite input tables displayed above give
\begin{center}
\begin{tabular}{r r r}
\toprule
$k$ & $g_k$ when available & $z_k$\\
\midrule
0&$1$&$0$\\
1&$-8$&$4$\\
2&$22$&$-15$\\
3&$-230/3$&$169/3$\\
4&$382/3$&$215/2$\\
5&$-22468/15$&$26449/15$\\
6&$-283418/45$&$232253/10$\\
7&---&$101570743/315$\\
\bottomrule
\end{tabular}
\end{center}
There is no missing input for $z_7$: the series $U^{-1}-1$ starts
in degree one, so $g_0,\ldots,g_6$ together with
$h_0,\ldots,h_7$ suffice. The seventh-order truncation therefore
has $O(n^{-8})$ error at fixed order.

\subsection{Scope of probabilistic transfer}
For any sequence of graph events $E_n$,
\[
 |\mathsf P_n(E_n)-\mathsf U_n(E_n)|\leq\frac4n+O(n^{-2}).
\]
For real graph statistics $Z_n$ with a uniform absolute bound $M$,
\[
 |\EE_{\mathsf P_n}Z_n-\EE_{\mathsf U_n}Z_n|
       \leq2M\,d_{\TV}(\mathsf P_n,\mathsf U_n)
       =\frac{8M}{n}+O_M(n^{-2}).
\]
If $0\leq Z_n\leq1$, the factor two can be omitted. An unbounded
statistic requires additional tail or uniform-integrability control.
The estimate is additive: an event of probability much smaller
than $1/n$ need not have similar relative probabilities under the
two laws. As with the enumeration series, no growing-order
truncation is justified by the all-orders statement.

\section{Source overlap and the boundary of the contribution}\label{sec:sources}

\subsection{Exact enumeration already available in 1990}
Bayoumi, El-Zahar, and Khamis~\cite{BEK} give prime-automorphism
information, P\'olya substitution, and a recursive generating-function
system. Their Lemmas 4.2 and 4.4 and Proposition 4.5 are prior art
for exact enumeration. Their Table 3, printed p.~4.47, agrees with
OEIS A123448 through order fourteen and continues through order
twenty. The original author-uploaded paper and the table page were
inspected; no asymptotic enumeration statement was located there.

For bibliographic completeness, the six table entries beyond the
inspected OEIS list are recorded below. They are \emph{quoted
prior-art table values}, not new enumerations certified by this
report's companion.
\begin{center}
\begin{tabular}{r r}
\toprule
$n$ & $a_n$ as printed in the 1990 table\\
\midrule
15&232072611578\\
16&3819032840830\\
17&66512264611298\\
18&1222446378826186\\
19&23650094948059080\\
20&480517922278457424\\
\bottomrule
\end{tabular}
\end{center}
The authors state that they used double precision. The remote table's
large-number typography and that numerical qualification are reasons
not to describe all its high-order entries as newly verified exact
integers. None of these six values enters the proof or the eight
reported asymptotic coefficients. Exact finite verification here is
bounded by order nine.

\subsection{Classical related asymptotics and sampling}
The publisher abstract of El-Zahar--Sauer~\cite{ES} states that the
number of pairwise nonisomorphic two-dimensional posets on $n$
elements is asymptotic to $n!/2$, using structural decomposition and
unique representation. This is a close classical antecedent to the
present leading equivalent. It counts posets rather than undirected
comparability graphs; nevertheless a first-ever claim for $n!/4$
would require inspection of the complete proof and its consequences.
The full subscription article was not obtained in this review.

Winkler's paper~\cite{Winkler}, dated December 1990 by the publisher
and listed as 1991 in an author bibliography, compares labeled,
unlabeled, and random-linear-order-pair models of dimension-two
orders. Its abstract reports the labeled equivalent
$(n!)^2/(2\sqrt e)$ and discusses rigidity, unique realizability,
and uniquely orientable comparability graphs. Its full subscription
text was likewise not inspected. These are sampling precedents,
so \eqref{eq:tvintro} is presented as a concrete sharp quantitative
refinement for the two explicitly defined graph laws, not as the
first comparison of related random-order models.

Literal transitive orientations on labeled vertices must be
separated from orientations modulo graph automorphisms. For example,
flipping twin-edge vertices can change a labeled orientation while
preserving an unlabeled poset type. The exact weighted identity in
Section~\ref{sec:four} avoids conflating these notions.

\subsection{Later permutation and intersection graph sources}
Johnston's 2020 account~\cite{Johnston} reports graph enumeration
through order fourteen and says its author did not know whether
the count was $o(n!)$. That is a historical statement about that
account, not a comprehensive assessment of the literature.

The inspected arXiv:2402.06394v1 (2024) of
Bassino--Bouvel--F\'eray--Gerin--Pierrot~\cite{BBFGP} studies
scaling and graphon limits and uses coarse bounds
$n!/30\leq a_n\leq n!$ in its permutation-graph discussion.
A separately available author-hosted PDF has February 2024 file
metadata and the same coarse lower bound; it is not evidence that
the final journal text has been inspected.

The final article, \emph{Dense and nondense limits for uniform
random intersection graphs}, appeared in \emph{Electronic Journal
of Probability} 31 (2026), article 110,
DOI~\href{https://doi.org/10.1214/26-EJP1570}{10.1214/26-EJP1570}.
The official landing and PDF requests returned an access-error
page in the source review. The final text remains uninspected, so
possible overlap with it is explicitly unresolved. No access-control
bypass, author contact, or purchase was used.

\subsection{What is established here}
The graph-specific result is the finite rooted-context and
large-prime-quotient approximation \eqref{eq:finite}, followed by
its all-orders coefficient formula and the stated inverse and
sampling consequences. Simple-permutation decomposition, its exact
functional equation, and its full factorial asymptotics are
established inputs. The current report and repository-index search
found no matching permutation-graph topic, but this was a bounded
absence check, not a universal novelty proof.

The conclusions that require no unresolved priority assertion are:
the proof gives the fixed-order expansion in the stated model;
the eight coefficients have independent exact checks; the inverse
windows preserve rounding; and the sharp total-variation constant
is four for the laws \eqref{eq:laws}. No OEIS submission, external
publication, or public repository modification is part of this
report.

\section{Reproducibility and verification limits}\label{sec:replay}

The source archive contains this LaTeX source, the rendered PDF,
standard-library Python exact-checking programs and tests,
machine-readable coefficient claims, a source-provenance record,
checksums, and deterministic release scripts. It contains neither
unrelated earlier reports nor private review documents. The root
README supplies the exact commands and expected outputs.

The fast companion replay checks rational series identities,
coefficient agreement by two routes, fixed-cutoff stabilization,
ordinary-power conversion, logarithmic coefficients, the symbolic
second-shift cancellation, and denominator divisibility. Exhaustive
graph replay is separately selectable: small orders provide a fast
regression test, while order nine reruns all finite inputs used for
the displayed eight coefficients. A full dual-route order-nine
replay is materially slower and is deliberately not hidden inside
the default tests.

Tests run under both ordinary Python and \texttt{python3 -O};
verification does not depend on the presence of Python
\texttt{assert} statements. Command-line outputs use new files or
new directories, reject existing destinations and symbolic-link
traversal, and avoid silently overwriting the report or a prior
result. Release manifests use a fixed allowlist, and the ZIP has
fixed member order, timestamps, and metadata. PDF rebuilding uses
fixed metadata, no shell escape, a clean TeX environment, and
reference passes to a stable auxiliary state. Byte-for-byte
reproduction is expected within the recorded toolchain; another
TeX or font version may typeset identically without producing
identical bytes.

The release checks certify the finite calculations and packaging,
not the analytic proof by executable means. In particular they do
not establish a numerical onset for the asymptotic inequalities,
turn the formal generating series into convergent analytic
functions, validate a growing truncation order, or resolve the
uninspected-source priority questions. Those boundaries are part
of the result rather than omissions that a finite successful test
can erase.

\section{Further research questions}\label{sec:future}

\paragraph{Effective remainder constants and threshold decisions.}
The interval-tail proof can be made numerical at each fixed order,
but an effective final error also requires explicit constants in
the simple-permutation transfer and in its finite-polynomial
composition. A useful next goal is to propagate all of these bounds
together, producing specified pairs $(K_R,n_R)$ for
\eqref{eq:logerror}. This would turn the asymptotic threshold window
into a certified large-input test: evaluate its two integer endpoints,
and, when they differ, settle the remaining case by an independent
exact bound on $a_n$. Such a certificate would need directed numerical
error control in the gamma and Lambert-$W$ evaluations as well as
the asymptotic remainder; floating-point agreement alone is
insufficient.

\paragraph{Coefficient growth and orders that increase with size.}
The rational-coefficient formal generator \eqref{eq:H} determines every fixed
coefficient, but the growth of $h_k$ is not bounded in this report.
Additional exact inputs could guide a study of coefficient growth,
while the proof would need estimates uniform in the cutoff $R$:
the number of bounded decorations, the interval-exclusion constant,
and the factorial-transfer remainder all change with $R$. Until
that dependence is controlled, neither a truncation order tending
to infinity nor an optimal truncation rule is justified. Questions
about Borel summation and exponentially small terms go further.
The symmetric-quotient estimate isolates one negligible scale, but
does not classify all sectors beyond the factorial Poincar\'e
expansion or their interactions. A complete transseries description
would require new information beyond the fixed-order theorem.

\paragraph{Efficient finite input recurrences.}
Exhaustive canonicalization makes the coefficient formula effective
at every finite order but rapidly becomes expensive. The 1990
P\'olya recurrence offers an established starting point for a faster
unrooted calculation. Extending exact recursive enumeration to
vertex-rooted contexts and, separately, to one-realizer graphs and
singleton root orbits could make many more $h_k$, $g_k$, and $z_k$
accessible. The marked multiplicity \eqref{eq:t} is the quantity
that must be retained: ordinary vertex pointing by $na_n$ does
not supply $r_n$, and dropping nontrivial marked contexts fails
already for seven-vertex contexts. A practical recurrence should first be
reconciled with both exhaustive routes on their overlapping range.

\paragraph{Finer fiber laws and sampling comparisons.}
Lemma~\ref{lem:fiberproduct} keeps more information than the event
$f(G)=4$: it factors the full fiber through a marked context and
all port fibers. One can therefore investigate other fixed fiber
sizes or bounded functions of fiber size by retaining these factors
rather than only selecting those equal to one. The principal
technical questions are how to organize the resulting finite input
data and how to control the weighted exceptional classes. For
unbounded positive weights, the existing class-count error is not
enough. Likewise, relative comparison of rare events or expectations
of unbounded graph statistics needs event-specific or tail estimates,
not merely the total-variation series proved here.

\begin{thebibliography}{99}
\bibitem{AAK}
M.~H. Albert, M.~D. Atkinson, and M.~Klazar,
\emph{The enumeration of simple permutations},
Journal of Integer Sequences \textbf{6} (2003), Article 03.4.4.
\href{https://cs.uwaterloo.ca/journals/JIS/VOL6/Albert/albert.html}{Journal article and full text}.

\bibitem{BEK}
B.~I. Bayoumi, M.~H. El-Zahar, and S.~M. Khamis,
\emph{Counting Permutation Graphs},
25th Annual Conference on Statistics, Computer Science and
Operations Research, Cairo University, 8--10 December 1990,
pp.~4.23--4.49.
\href{https://www.research.asu.edu.eg/handle/123456789/1142}{Institutional record};
\href{https://www.researchgate.net/profile/Soheir-M-Khamis/publication/324587746_Counting_Permutation_Graphs/links/5ad7259ba6fdcc2935835c6e/Counting-Permutation-Graphs.pdf}{author-uploaded full text}.

\bibitem{Borinsky}
M.~Borinsky,
\emph{Generating asymptotics for factorially divergent sequences},
Electronic Journal of Combinatorics \textbf{25}(4) (2018), P4.1.
\href{https://www.combinatorics.org/ojs/index.php/eljc/article/download/v25i4p1/pdf/}{Published full text}.
Theorem 35 and equations (62)--(67).

\bibitem{HP}
M.~Habib and C.~Paul,
\emph{A survey of the algorithmic aspects of modular decomposition},
Computer Science Review \textbf{4}(1) (2010), 41--59.
\href{https://www.irif.fr/~habib/Documents/HP10.pdf}{Author-hosted final text}.
Lemma 20, printed p.~55.

\bibitem{BBFGP}
F.~Bassino, M.~Bouvel, V.~F\'eray, L.~Gerin, and A.~Pierrot,
\emph{Dense and nondense limits for uniform random intersection
graphs}, arXiv:2402.06394v1 (2024).
\href{https://arxiv.org/abs/2402.06394v1}{Exact inspected version}.
The associated final article is
\emph{Dense and nondense limits for uniform random intersection
graphs}, Electronic Journal of Probability \textbf{31} (2026),
article 110, 44 pp.
\href{https://doi.org/10.1214/26-EJP1570}{Final journal record}.
The final text was not inspected.

\bibitem{ES}
M.~El-Zahar and N.~W. Sauer,
\emph{Asymptotic enumeration of two-dimensional posets},
Order \textbf{5} (1988), 239--244.
\href{https://doi.org/10.1007/BF00354891}{Publisher abstract and record}.
Full subscription text not inspected.

\bibitem{Winkler}
P.~Winkler,
\emph{Random orders of dimension 2},
Order \textbf{7}, 329--339 (publisher date December 1990;
author bibliography 1991).
\href{https://doi.org/10.1007/BF00383197}{Publisher abstract and record}.
Full subscription text not inspected.

\bibitem{OEIS}
The OEIS Foundation,
\emph{A123448 Number of permutation graphs on n vertices}.
\href{https://oeis.org/A123448}{OEIS entry}.
Inspected source snapshot, 3 October 2026.

\bibitem{Johnston}
T.~Johnston,
\emph{Enumerating permutation graphs}, 25 October 2020.
\href{https://tomjohnston.co.uk/blog/2020-10-25-enumerating-permutation-graphs.html}{Author's enumeration account}.
\end{thebibliography}
\end{document}
